\section{Context and further questions}\label{sec:outlook}
\subsection{What is inherited and what is specialized}
The leading equivalent is established prior work of Canfield--McKay \cite{CM} and is also within Barvinok--Hartigan's smooth-margin theory \cite{BH}. General complex cumulant control is due to Isaev \cite{Isaev}, with earlier complex-martingale antecedents in Isaev--McKay \cite{IM2018}. Connected-Wick identities and high-dimensional power counting are standard tools, explicitly developed in the related enumeration literature \cite{IM,IMZ}. Higher-order saddlepoint corrections for contingency tables also have a substantial earlier history, including the double-saddlepoint approach of Zipunnikov--Booth--Yoshida \cite{ZBY}.

The assertion proved here is the stated geometric square-table specialization with its complete cutoff and gauge reduction, all-fixed-order rational coefficient algorithm, evaluated corrections, and controlled consequences. It is not a claim to a new general saddlepoint or cumulant method. A bounded primary-source review did not locate the explicit coefficient pair or this dense-square specialization, but absence from the sources inspected is not a worldwide priority certificate. The sparse occurrence of the first numerical coefficient is noted in Section~\ref{sec:conjecture}.

A concrete comparison explains why a leading-order Edgeworth expression does not settle the present correction. Specializing the cubic/quartic exponent displayed by Barvinok--Hartigan to these square margins gives, by direct algebra,
\begin{equation}\label{eq:BHcompare}
-\frac\mu2+\nu=B_0(u)-\frac1n+
\frac1{n^2}\left(\frac7{12}+\frac1{12u}\right).
\end{equation}
At $\lambda=1$, this is $1/4-1/n+5/(8n^2)$. Their theorem uses this expression with relative error tending to zero; it does not assert that its lower-order displayed terms are the true higher-order counting coefficients. The additional connected cumulants account for the difference at the precision sought here. Equation~\eqref{eq:BHcompare} is a specialization of their formula, not a correction to a claim in their theorem.

Recent binary or bounded-entry enumeration results are likewise not interchangeable with geometric tables. For example, admixed arrays in Aw \cite{Aw} and simple multipartite hypergraphs in Isaev--Makai--McKay \cite{IMM} concern different underlying models. Their statements do not supply the analytic specialization above.

\subsection{Questions opened by the expansion}
\begin{enumerate}
\item \textbf{Explicit finite-size bounds.} Can the localization and cumulant constants be tracked sharply enough to certify a useful numerical remainder and a starting index? This would turn the inverse two-ceiling statement and the eventual correction-factor inequality into finite-size certificates.
\item \textbf{Higher coefficients and efficient enumeration.} The finite algorithm proves computability but grows rapidly. Structural reductions of colored connected graphs, and symbolic use of the density polynomials, may make several additional orders practical. The independent moment-factorization route offers a useful second implementation.
\item \textbf{Growth and optimal truncation.} What is the growth of $P_j$ with $j$? Fixed-order proofs do not answer whether the series converges, has a Gevrey bound, or admits useful optimal truncation. Identifying exponentially small contributions would require analysis beyond the present Poincar\'e expansion.
\item \textbf{Transition regimes.} The compact-density constants degenerate as density tends to zero or infinity. A parameter-dependent analysis could connect the positive-density result with sparse table asymptotics or growing-margin regimes. Substitution into a fixed-density formula is insufficient to justify those limits.
\item \textbf{Rectangular and nonconstant margins.} The null direction remains, but the whitening map, covariance structure and coefficient complexity change. Extending the full argument to controlled aspect ratios or smooth nonconstant margins requires fresh uniform bounds rather than an automatic appeal to the square proof.
\item \textbf{The correction-factor conjecture at all sizes.} Corollary~\ref{cor:delta} identifies a positive eventual correction in a dense square regime. Proving the conjectured strict inequalities at every finite size, or uniformly across changing densities and aspect ratios, remains a different problem.
\end{enumerate}
These questions preserve the distinction between the theorem already proved, its computational checks, and stronger conclusions that would need additional analysis.

The theorem numbering for Canfield--McKay and the complex cumulant result refers to the linked preprint versions. The source package links to primary papers but does not redistribute their PDFs.
